\section{Finite descriptions and stopped numerical certificates}
\label{sec:numerics}
An arbitrary Borel acceptance set is not a finite computer program. We first state the error bound on the scale that a stopped experiment actually uses.

\begin{proposition}[Relative-acceptance implementation]\label{prop:relative}
Consider an implemented attempt program with the same legal physical ports and $B$ live states as the ideal one. On every surviving pre-attempt history suppose its conditional expected cost is at most $\widetilde\kappa$, its acceptance probability is at least $(1-\rho)/b$, and its unnormalized terminal loss is at most $(Q+\eta)/b$, where $0\le\rho<1$, $\eta\ge0$, $b\ge1$. It then terminates almost surely and satisfies
\begin{equation}\label{eq:relative}
 \E C\le\frac{\widetilde\kappa b}{1-\rho},\qquad
 R\le\frac{Q+\eta}{1-\rho}.
\end{equation}
The terminal alphabet is the actual exported dictionary. These bounds remain valid for unknown history-dependent implementation errors satisfying the stated conditional certificates.
\end{proposition}
\begin{proof}
The probability of surviving $j$ attempts is at most $(1-(1-\rho)/b)^j$. Predictable summation of the bounded conditional costs and nonnegative terminal-loss contributions proves both inequalities and finite expected attempts. Each attempt terminates almost surely by its own cost certificate. No comparison of terminal distributions was used to infer duration.
\end{proof}

For example, an absolute one-attempt acceptance error $e$ and unnormalized loss error $e_L$ give $\rho=be$ and $\eta=be_L$, not $\rho=e$ and $\eta=e_L$. Rare selection amplifies these absolute errors. A task-readout perturbation of size $\delta$ in a carrier of diameter $D_H$ changes loss by at most $2D_H\delta+\delta^2$ \emph{conditional on acceptance}; its unnormalized error is weighted by acceptance mass. A continuous report and a digitized atomic report need not be close in total variation; boundary-event or unnormalized moment estimates must be supplied instead.

\begin{proposition}[Finite-description upper at the same cost]\label{prop:finite}
In addition to exact balance, assume the known reference measure admits an effective sequence of finite physical report partitions with certified approximation of bounded acceptance functions in $L^1(\nu)$. Assume finite points in $H$ and their exported values can be approximated within $H$ with certified uniform error, and their primitive report classifiers and probabilities have declared finite descriptions, execution costs and integration error bounds. Then, for every $b\ge1$, $m\ge1$ and $e>0$, there is a finite-description observer with at most $B+D_0m$ states, cost at most $\kappa b$ and risk at most $\Qm_m(b)+e$. This is an existence and approximation result under effective certificates, not a polynomial-time discovery theorem for arbitrary kernels.
\end{proposition}
\begin{proof}
For $b>1$, choose $1<b'<b$. By Lemma~\ref{lem:attain}, $\Qm_m(b')$ can be made arbitrarily close to $\Qm_m(b)$. Approximate its optimal finite dictionary inside $H$ and use its nearest-center cells. Approximate the bounded acceptance $h/b'$ and this finite center assignment jointly on a finite partition. The assumptions supply certified $L^1(\nu)$ error bounds. Because all squared losses are bounded and dictionary errors change them uniformly, the acceptance and unnormalized loss errors tend to zero. The ideal acceptance $1/b'$ has a strictly positive margin above $1/b$. Take the approximation sufficiently accurate that its acceptance is still at least $1/b$. Balance makes these quantities the same on every surviving history and model. Its risk is its unnormalized loss divided by acceptance; it tends to the ideal risk. The attempt count proof then gives cost at most $\kappa b$. For $b=1$, accept the first report identically and approximate only the dictionary and assignment; no acceptance margin is needed. Current partition indices and command coins are scratch variables, and only the routing or exported output label persists.

The argument is constructive once the indicated approximation and integration certificates are supplied. A finite search over candidate dictionaries and report partitions, with bounded integration intervals, can certify an available upper. No claim that a local optimizer supplies a global lower is involved. Fully effective global convergence additionally uses effective compactness of the dictionary carrier and uniform integration bounds; a finite dictionary net and the Lipschitz loss then give finite lower and upper profile approximations, potentially at exponential search cost.
\end{proof}

For regular geometry, finitely supplied centers and coordinate certificates give a direct upper without learning an arbitrary acceptance density. Replace inner rectangles from Proposition~\ref{prop:geometry} by outward rounded rectangles, accepting every true inner rectangle and only points with marked distance at most a fixed multiple of its radius. Actual mass remains at least $1/b$ and distortion remains of the same order. With $m$ centers in a fixed $k$-dimensional box and radius $t$, $p=O(1+\log(1/t))$ precision bits suffice when the coordinate evaluation has that error certificate. A sequential implementation uses $O(mkp)$ read-only data bits, $O(kp+\log m)$ scratch bits and $O(mkp^2)$ elementary arithmetic operations per productive report, in addition to the specified physical evaluation and command costs. These are sufficient costs, not optimal bit/time lower bounds.

\subsection*{An end-to-end dyadic example at zero positivity floor}
Use the classical family of Theorem~\ref{thm:classical} with independent uniform $(Y,T)\in[0,1]^2$, post-command audit $Z=Y+TE$, and $p_0=1/4$. Thus $\kappa=5$. Use three states: gate, acquisition, and terminal value $1/2$. Ideally accept
\[
 A=[3/8,5/8]\times[0,1/8].
\]
Under balanced task measures,
\begin{equation}\label{eq:dyadic-ideal}
 \nu(A)=1/32,\qquad
 \int_A[(y-1/2)^2+t^2]\,dy\,dt=1/3072.
\end{equation}
Let $e=2^{-14}$, and suppose the physical digitizer supplies clipped coordinates with errors at most $e$, using only $(Y,T)$ and not the discarded sign-flip coin. Accept when
\[
 3/8+e\le\widehat Y\le5/8-e,\qquad \widehat T\le1/8-e.
\]
The accepted true report lies in $A$, and every point in
$A_{\rm in}=[3/8+2e,5/8-2e]\times[0,1/8-2e]$ is accepted. Put
\begin{equation}\label{eq:dyadic-constants}
 \alpha_-=(1/4-4e)(1/8-2e),\quad
 \beta_+=1/3072,\quad
 W_G=5/\alpha_-,\quad W_A=W_G-4,\quad
 U_G=U_A=\beta_+/\alpha_-.
\end{equation}
These are exact rational numbers. The reference rows with gate success $1/4$ and productive acceptance $\alpha_-$ obey
$W\ge\one+QW$ and $U\ge\beta+QU$. Every actual gate row has greater success and every implemented productive row has acceptance at least $\alpha_-$ and loss contribution at most $\beta_+$. Thus the same inequalities hold directly for the implemented raw rows: a killed-potential certificate has relative perturbation constants $\rho=\eta=\zeta=0$ relative to these conservative rows. Summing the finite-horizon inequalities, \emph{including} the nonnegative surviving $W$ and $U$ terms before taking the limit, proves
\begin{equation}\label{eq:dyadic-final}
 \E C\le 5/\alpha_-<161,\qquad R\le1/(3072\alpha_-).
\end{equation}
Equivalently these bounds follow from Proposition~\ref{prop:relative} with the conservative acceptance and numerator. A same-report changed decision lies only in $A\setminus A_{\rm in}$, whose probability is at most $1/32-\alpha_-$. No total-variation comparison of continuous and atomic report laws was made.

For a matching all-program lower at the same budget $161$ and three states, Theorem~\ref{thm:main} gives $\Qm_1(161/5)$. For any output $r$, a sublevel $t^2+(y-r)^2\le z$ in $[0,1]^2$ has area at most $\pi z/2$. Integration of the selected tail gives $\Qm_1(b)\ge1/(\pi b)$. When $b\ge8/\pi$, a half-disk centered at $(1/2,0)$ of area $1/b$ stays inside the square and attains equality. Consequently
\begin{equation}\label{eq:dyadic-interval}
 \frac5{161\pi}=R_\aut^\Theta(161,3)
 \le R_{\rm dyadic}\le\frac1{3072\alpha_-}.
\end{equation}
The equality is the ideal Borel optimum; the displayed dyadic observer is a separately certified nearby finite implementation.

The three threshold integers are $6145,10239,2047$, requiring 45 read-only bits at 15 bits each, plus the fixed comparator program and output $1/2$. Two 15-bit digitized-coordinate registers, one fresh flip bit and fixed control scratch suffice within a productive call. The three persistent states require two encoding bits; neither a counter nor a retained random seed is present. Comparisons use $O(14)$ bit operations per productive call. If a gate call costs at most $t_G$ time and productive preparation, digitization and comparison together cost at most $t_A(14)$, expected physical execution time is at most $161\max(t_G,t_A(14))$. The promised digitizer and balanced physical pulse costs are not free. One subsequent physical audit makes the total expected raw-call budget at most $162$.
